% ====================================================================
% s31_map.tex — §3.1 생존 지도 본문화 (W2 초안 — 실측·데이터 강조축)
%   원천 = ch1_appD_si(부록 예비 지도)를 장 본문체로 승격. 노드 판정(이월/재해석/새물리/부분/구조이월)은
%   그대로 보존(P5) — 부록체→본문체·Part 0/I 결과 xr live 인용. 라벨은 신규 sec:si-*/tab:si-*.
%   (마스터 주: 본 절이 ch1_appD_si 를 대체하면 ch3v22_notation 의 \S\ref{ssec:si-gap} 는
%    본 패키지 §3.4 sec:si-mech 로 재지정 필요 — DESIGN_NOTE 라벨 이관표 참조.)
% ====================================================================
\section{Si 계열 접목의 생존 지도 --- 무엇이 이월되고 무엇이 새 물리인가}\label{sec:si-survival}
Chapter 1$\cdot$2 는 흑연 음극과 LCO 양극에서 같은 골격(전하 보존 반전~\eqref{eq:sm-mc-balance} 을 중심식으로
하는 평형 종$\to$합산 사슬)이 전극을 가리지 않고 성립함을 보였다. 세 번째 활물질 계열인 실리콘은 삽입이
아니라 \emph{합금화}로 리튬을 저장하므로 앞 두 장과 물리가 갈라지는 지점이 많다. 본 절은 완결 유도를
두지 않고 --- 골격 노드 $\mathrm{N0}$--$\mathrm{N9}$ 를 실리콘에 대조해, 식$\cdot$부호가 그대로 사는
곳$\cdot$형식은 유지하되 변수 의미가 바뀌는 곳$\cdot$골격 밖 새 물리가 필요한 곳을 가르는 \emph{지도}를 먼저
놓는다. 지도의 전 서술은 문헌 기반이며(자체 실측 대조 없음), ``새 물리'' 로 분류된 것의 취급은
\S\ref{sec:si-blend}(전하 보존의 블렌드 일반화)$\cdot$\S\ref{sec:si-mech}(기계 히스테리시스)로 넘긴다.

\subsection{Si 가 흑연과 다른 곳 --- 확립 사실 여덟}\label{ssec:si-eight}
지도의 근거가 되는 실측 사실 여덟을 먼저 고정한다. (i) Si 는 삽입이 아니라 \emph{합금화}다 --- 고온 평형은
네 중간상(Li$_{12}$Si$_7$ 등)의 계단 plateau 를 보이나\cite{wen_huggins1981} 호스트 격자 자체가
재구성된다. (ii) 상온 첫 리튬화는 결정질 Si 가 비정질 Li$_x$Si 로 바뀌는 \emph{전기화학적 고상
비정질화}이며, 입자 안에서 결정 코어$/$비정질 껍질의 날카로운 이동 경계로 진행한다(입자 내
두-상)\cite{limthongkul2003,li_dahn2007}, 그 코어--쉘 이동은 in~situ TEM 으로 직접 관측된다\cite{mcdowell_coreshell2013}.
(iii) 깊은 리튬화 끝(약 $50$ mV)에서 준안정 결정상 Li$_{15}$Si$_4$ 가 갑자기 결정화하며\cite{obrovac_christensen2004},
이 상전이는 operando $^7$Li NMR 로 확증된다\cite{ogata_nmr2014}. (iv) 첫 사이클 이후의 비정질 경로는
plateau 없는 \emph{완만한 경사 전위} $U(x)$ 이고, 제일원리 계산이 이 곡선을 재현한다\cite{chevrier_dahn2009}
--- 나노 a-Si 는 첫 사이클만 두-상이고 이후 순환은 solid-solution 경사로 옮겨간다\cite{wang_asi2013}. (v)
만충 부피 변화는 약 $300\%$ 로\cite{beaulieu2001} 흑연($\sim10\%$)과 자릿수가 다르다. (vi)
충$\cdot$방전 갈림(히스테리시스)은 수백 mV 급이고, in~situ 응력 측정이 소성 유동(압축 $\sim-1.75$
GPa)과 \emph{응력--전위 결합} $\sim100$--$120$ mV/GPa 를 직접 보인다\cite{sethuraman_stressevo2010,
sethuraman_stresspot2010} --- 히스테리시스의 지배 몫이 \emph{기계적}이다. (vii) 임계 입자경 $\sim150$ nm
아래에서만 첫 리튬화 균열이 없다 --- 입자 크기가 1차 인자다\cite{liu_sizefracture2012}. (viii) 이 전부의
총람은 합금 음극 리뷰\cite{obrovac_chevrier2014} 와 나노 Si 설계 리뷰\cite{mcdowell_review2013} 가 준다.
\begin{srcbox}
근거 등급 --- 위 여덟은 모두 1차 문헌 실측(사실 vi 응력 결합의 쪽 범위는 사용 시 Crossref 대조 유보,
tier A). 사실 ii$\cdot$iii 의 관측 채널(XRD$\cdot$TEM$\cdot$NMR)이 상보적으로 겹치는 두-상 전이 둘이
\S\ref{ssec:si-demo} 의 골격 시연 자리다. 정량 수치의 케이스별 정리는 \S\ref{sec:si-cases}
표~\ref{tab:si-cases} 가 담당한다.
\end{srcbox}

\subsection{노드 생존 판정표}\label{ssec:si-nodemap}
판정 다섯 갈래: \textbf{이월}(식$\cdot$부호 그대로) $/$ \textbf{구조 이월}(구조는 살고 미시 해석만 재작성) $/$
\textbf{재해석}(형식 유지, 변수 의미 변경) $/$ \textbf{부분}(일부 조건에서만) $/$ \textbf{새 물리}(골격에 없는 항
필요). 표~\ref{tab:si-nodemap} 이 흑연으로 세운 골격 노드의 실리콘 생존 판정이며, 판정 열은 Chapter 1
부록 예비 지도의 결론을 본문 결과 사슬 순서로 옮긴 것이다(판정 자체는 불변).

\begin{table}[h]
\centering\small
\begin{tabular}{p{0.12\textwidth} p{0.13\textwidth} p{0.63\textwidth}}
\toprule
노드 & 판정 & 무엇이 살고 무엇이 깨지나 \\
\midrule
N0 조건 & \textbf{이월} & $\sigma_d$$\cdot$$|I|$ 환산은 전극 무관 --- 저전위 음극 규약 그대로. \\
N1 분극 & 재해석 & 옴 분극 벗기기는 성립하나, $|I|\to0$ 에서도 안 사라지는 \emph{기계(응력) 과전위 몫}이 $V_n$ 에 남는다(사실 vi) --- \S\ref{sec:si-mech}. \\
N2 중심 & 재해석 & 이산 staging 중심 $U_j$ 부재(사실 iv) --- 유효 logistic 성분의 파라미터화(MSMR-Si\cite{verbrugge_lisi2016}) 또는 연속 $U(x)$\cite{chevrier_dahn2009} 로 대체. $\partial U/\partial T=\Delta S/F$ 관계 자체는 생존. \\
N3 히스 & \textbf{새 물리} & 정규용액 spinodal gap 상한($\Omega{=}4RT$ 에서 $\approx55$ mV, 그림~\ref{fig:hysgap})으로 수백 mV 를 못 담는다 --- \S\ref{sec:si-mech}(GS-1). \\
N4 폭 & 재해석 & 경사 성분의 ``폭''은 합금 조성폭(수십 \%)이지 열적 $RT/F$($\sim26$ mV)가 아님 --- $n_j\gg1$ 의 현상학 폭. \\
N5 $\xi_\eq$ & 구조 이월 & logistic 합 \emph{형식}은 생존(detailed balance 는 전극 무관; MSMR-Si 실증\cite{verbrugge_lisi2016}) --- ``동등 격자 자리'' 미시 해석은 비정질에서 소멸, 유효 파라미터화로 재해석. \\
N6 peak & 부분 & 날카로운 peak 문법은 Li$_{15}$Si$_4$ 탈리튬화$\cdot$1차 리튬화 두-상(사실 ii$\cdot$iii)에만 --- 나머지는 넓은 hump. \\
N7 broadening & 부분 & 두-상 broadening 틀은 부분 적용되나 기작이 다르다(입자 간 순차가 아니라 \emph{입자 내} 코어--쉘\cite{mcdowell_coreshell2013}). 흑연에서 정량 배제한 입자-크기 채널이 Si 에선 1차 인자로 역전(사실 vii). \\
N8 지연 & 구조 이월 & Eyring 유한율속 꼬리 구조는 전극 무관. 단 $|I|\to0$ 에서 꼬리는 소멸해도 기계 히스는 남으므로, ``잔여 갈림 $=$ 평형 분기'' 분리 논법이 약해진다. \\
N9 합산 & \textbf{이월} & 전하 보존 합산~\eqref{eq:sum} 은 전극 중립 --- 블렌드 일반화의 앵커, \S\ref{ssec:si-neutral}$\cdot$\S\ref{sec:si-blend}. \\
\bottomrule
\end{tabular}
\caption{골격 노드의 Si 생존 판정 --- 이월 2(N0$\cdot$N9)$\cdot$구조 이월 2(N5$\cdot$N8)$\cdot$재해석
3(N1$\cdot$N2$\cdot$N4)$\cdot$부분 2(N6$\cdot$N7)$\cdot$새 물리 1(N3). 근거 문헌은 \S\ref{ssec:si-eight}.
판정 분포는 Chapter 1 부록 예비 지도와 동일하다(본문 승격 시 판정 불변 --- 역사적 인계 보존).}
\label{tab:si-nodemap}
\end{table}

\subsection{살아남는 최강 자산 --- 전하 보존의 전극 중립성}\label{ssec:si-neutral}
본문 중심식 --- 전하 보존 음함수~\eqref{eq:sm-mc-balance} --- 는 격자 통계가 아니라 \emph{전하
회계}에서 왔으므로(\S\ref{sec:sm-mc}: 반응 몫의 용량가중 합 $=$ 통과 전하), staging 이 무너져도
살아남는다. 합금 전극의 다중 전기화학 반응$\cdot$화학종 분화(speciation)를 정확히 이 구조 --- 반응별
logistic 점유의 용량가중 합을 고정 전하에서 뒤집기 --- 로 세운 정식화가 Li--Si 에서
실증되어 있고\cite{verbrugge_lisi2016}, 이는 본문 MSMR 대응(\S\ref{sec:lco-code})과 같은 저자 계보의 Si
판이다. 곧 \emph{내부 전위를 결정하는 중심 구조는 흑연$\cdot$LCO$\cdot$Si 에 공통}이며, Si 접목에서 바뀌는
것은 이 구조에 들어가는 성분(전이 집합$\cdot$폭$\cdot$히스 항)이지 구조 자체가 아니다. 이 전극 중립성이
\S\ref{sec:si-blend} 의 블렌드 일반화 --- 두 host 가 \emph{하나의} 전하 보존식을 공유하는 확장 --- 을
가능케 하는 앵커다.

\subsection{부분 적용 시연의 자리 --- 진짜 이산 전이 둘}\label{ssec:si-demo}
골격 문법이 Si 에서 그대로 시연 가능한 곳은 진짜 이산 전이 둘뿐이다: (i) \textbf{1차 리튬화 두-상}
(c-Si$\to$a-Li$_x$Si 이동 경계 --- 두-상 delta$\to$종 broadening 문법 부분 적용, 단 입자 내 코어--쉘
기하로 재해석\cite{mcdowell_coreshell2013}), (ii) \textbf{Li$_{15}$Si$_4$ 결정화$\cdot$탈리튬화}(날카로운
dQ/dV 특징 --- 평형 종 문법 적용, NMR 확증\cite{ogata_nmr2014}). 나머지 경사 영역은 유효 다성분
logistic(MSMR-Si)으로만 파라미터화되며, 그 성분들은 물리 전이가 아니라 곡선 기술 성분이다 ---
흑연$\cdot$LCO 에서 성분$=$상전이였던 대응이 Si 에서는 깨진다는 것을 성분 해석에 명시해야 한다(이 재해석은
\S\ref{sec:si-cases} 의 케이스별 개형이 존중한다). 위 두 시연의 실제 수치 완주(초기값 표$\cdot$계산
예제)는 Si 실측 dQ/dV 확보 후의 소관이며, 본 장은 검토용 초기값 표(\S\ref{sec:si-cases})까지 놓는다.

\begin{keybox}
\textbf{생존 지도 한 줄.} 골격에서 \emph{전하 보존 중심식은 이월}(N9 --- 블렌드 앵커), \emph{평형 종
형식은 구조 이월}(N5 --- 성분 의미만 재해석), \emph{기계 히스테리시스만 새 물리}(N3 --- \S\ref{sec:si-mech}).
곧 Si 는 골격을 버리는 것이 아니라 성분을 갈아 끼우는 활물질이며, 유일한 골격 밖 항이 응력--화학퍼텐셜
결합이다.
\end{keybox}

% 자산(v1.0.22 R5 신규): [V22-R5-02 Si 사실 8건 문헌 지도 본문화 ssec:si-eight(NMR·TEM 관측채널 실측 앵커 추가)]
%   [V22-R5-03 노드 생존 판정표 본문판 tab:si-nodemap(판정 불변·xr live 결과 인용)]
%   [V22-R5-04 전하 보존 전극 중립성=블렌드 앵커 ssec:si-neutral] [V22-R5-05 부분 시연 자리 ssec:si-demo(성분≠상전이)]
